\section{The Dex programming model and its core semantics}
\label{sec:dex}

\subsection{The model}
\label{sec:dex-model}

Table~\ref{tab:constructs} lists the constructs of the Dex programming
model. The specification gives rules for all of them. The theorem needs only
those in the last column, because the translation of \S\ref{sec:translation}
uses no others.

\begin{table}[t]
\centering
\small
\begin{tabularx}{\textwidth}{@{}P{2.5cm}LP{2.6cm}@{}}
\toprule
Construct & What it is & Used by the translation \\
\midrule
Flow & One run of a Flow definition: the unit of execution and of
  observation. It starts with a start Step and an input, and it closes with
  a status and a result. & Yes \\
Step & A Step execution has an input, an optional \kw{WaitFor} method, and a
  required \kw{Execute} method. These two methods are the only place where
  application code runs, apart from RPCs. & Yes \\
StepDecision & What \kw{Execute} returns: \kw{GoTo}/\kw{GoToMany} (create
  Step executions), \kw{DeadEnd} (end this branch), \kw{ForceComplete} and
  \kw{ForceFail} (close the Flow), \kw{GracefulComplete},
  \kw{ForceCompleteIfChannelsEmpty}, \kw{CancelSteps},
  \kw{CancelSiblingSteps}. & \kw{GoTo}, \kw{GoToMany}, \kw{DeadEnd},
  \kw{ForceComplete}, \kw{ForceFail} \\
Wait & What \kw{WaitFor} returns: Timer, Channel, and SubFlow conditions
  combined by \kw{AllOf}, \kw{AnyOf}, or \kw{AnyCombinationOf}; or the empty
  wait. & \kw{AnyOf} and \kw{AllOf} over Timers and \kw{ForOne} Channel
  conditions \\
Attribute & A key-value store of one run. An AttributeMap is a family of
  Attributes; a Step method sees an AttributeMap instance only if its
  options select it. & Yes \\
Channel & A FIFO queue of one run. A Channel condition has bounds; the
  default, \kw{ForOne}, waits for one message and consumes exactly one. &
  \kw{ForOne} only \\
Timer & A wait condition with a duration in whole seconds. & Yes \\
SubFlow & A child Flow started by a wait; its result is delivered to the
  parent. & No \\
RPC & A client call that runs a handler, reads and writes the run, and may
  start or cancel Steps. & No \\
Attribute lock & Mutual exclusion on Attribute keys for the duration of one
  Step method or RPC. & No \\
\bottomrule
\end{tabularx}
\caption{The constructs of the Dex core. The specification covers all of
them; the Lean model has the ones the translation and the admissible
environment use.}
\label{tab:constructs}
\end{table}

Two properties of the model shape everything that follows. First,
application code never blocks. A Step that must wait says so by returning a
wait from \kw{WaitFor}; the engine runs \kw{Execute} when the wait is
satisfied. Second, application code does not touch shared state directly. A
method receives a \emph{snapshot} of the Attributes when it starts, and it
returns its writes, its publications, and its decision as one response.
The engine commits that response later, in one atomic step, unless the
Step was canceled in between; then the response is discarded.

\subsection{Configurations}
\label{sec:dex-config}

The specification defines the system state as partition clocks, a program,
the runs of all Flows, engine state, and cross-Flow messages in transit. The
admissible environment of the theorem starts one Flow, and the translation
uses no SubFlows, so the Lean model (\file{Dex/Config.lean}) reduces the
state to one partition:
\[
\Sigma \;=\; \langle\, \mathit{clock},\ \mathit{lastCommitAt},\
\mathit{run} : \mathit{Run}^{?},\ \Theta : [\mathit{Task}],\
\mathit{nextTaskId} \,\rangle .
\]
A run holds a status (\emph{Running} or \emph{Closed} with a close status),
the collected outputs, the Step executions, a per-type Step counter, the
Attributes, the pending messages of each Channel in commit order, the armed
timers with their fire times, and a sequence number $\mathit{seq}$ that
every commit increments. A Step execution has an ID $T\text{-}n$ (its type
and its number from the per-type counter), an input, options, a phase
$\varphi \in \{\mathit{WF}, \mathit{WT}, \mathit{EX}, \mathit{CO},
\mathit{CA}\}$ (WaitFor, Waiting, Execute, Completed, Canceled), an optional
\emph{invocation} (a method that is pending or running), an optional
registered wait, and optional condition results. A registered wait records
$\mathit{sws}$, the sequence number at which it was registered; it orders
waiters.

$\Theta$ is engine state that the persistent state machine does not hold:
the \emph{tasks}, which are method attempts handed to a Worker. A task
carries the request the engine built when the method started (context,
input, snapshot, condition results) and a phase: \emph{Started},
\emph{Running} with the effects still to emit and the response, or
\emph{Returned} with the response.

\subsection{Rules}
\label{sec:dex-rules}

A transition is $\Sigma \xrightarrow{\;\vec{\ell}\;} \Sigma'$, where
$\vec{\ell}$ is the ordered sequence of observable events it produces. A
rule marked \emph{commit} is one command of the server's deterministic
state machine: it is atomic, it stamps the current clock value, and the
clock must have strictly increased since the previous commit. In Lean a rule
is a case of the partial function
\lean{Dex.step} (\file{Dex/Rules.lean}): the \emph{action} names the rule and
carries its nondeterministic choice, so determinism of each rule holds by
construction (\lean{step_deterministic}). Table~\ref{tab:rules} lists the
ten actions.

\begin{table}[t]
\centering
\small
\begin{tabularx}{\textwidth}{@{}P{3.1cm}P{1.9cm}LP{2.9cm}@{}}
\toprule
Action (rule) & Kind & Effect & Events \\
\midrule
\lean{tick} \rulename{TICK} & environment & The clock advances by
  $\delta > 0$ microseconds. & none \\
\lean{request} \rulename{FLOW-START} & commit & Enabled when no run exists.
  Creates the run and its start Step. & $\ev{start}(T_0, v)$ \\
\lean{request} \rulename{PUBLISH} & commit & While the run is running:
  appends the message and wakes the Channel's waiters. Otherwise the request
  is rejected and nothing changes. & $\ev{pub}(c, v)$, or
  $\ev{rejected}(\ev{pub}(c, v))$ \\
\lean{stepStart} \rulename{STEP-START} & commit & A pending invocation
  becomes running. The snapshot is taken, and a fresh task enters
  $\Theta$. & none \\
\lean{workerRun}, \lean{workerEffect}, \lean{workerReturn}
  \rulename{WORKER-EFFECT} & Worker & The handler runs; the task emits its
  external effects one at a time; the task becomes \emph{Returned}. &
  $\ev{fx}(e, m, a)$ per effect \\
\lean{stepComplete} \rulename{STEP-COMPLETE-*} & commit & The response of a
  \emph{Returned} task of a still-running invocation commits: its data
  effects, then the wait registration (WaitFor) or the decision (Execute). &
  none, or $\ev{closed}(\ldots)$ \\
\lean{workerAbandon} \rulename{WORKER-ABANDON} & Worker & A task whose
  invocation was canceled is dropped, before or after any of its effects. &
  none \\
\lean{stepCompleteStale} \rulename{STEP-COMPLETE-STALE} & Worker & A
  \emph{Returned} task whose invocation is no longer running is dropped. &
  none \\
\lean{timerFire} \rulename{TIMER-FIRE} & commit & Fires a non-empty prefix
  of the due timers, in index order, in one commit. & none, or
  $\ev{closed}(\ldots)$ \\
\bottomrule
\end{tabularx}
\caption{The transitions of the restricted Dex system
(\lean{Dex.Action}, \file{Dex/Rules.lean}). The specification also has rules
for the other client inputs, RPCs, SubFlows, and a deferred close; the
admissible environment and the translation never enable them.}
\label{tab:rules}
\end{table}

\paragraph{A Step method is not atomic.}
This is the central fact of the semantics. A Step method begins with a
Start commit (\rulename{STEP-START}). For an invocation that stays current
and returns, Start is followed by a sequence of Worker transitions and
then a Complete commit (\rulename{STEP-COMPLETE-WAITFOR} or
\rulename{STEP-COMPLETE-EXECUTE}). The Worker transitions are the parts of
\rulename{WORKER-EFFECT}, which the Lean model has as separate actions: one
\lean{workerRun}, zero or more \lean{workerEffect}, and one
\lean{workerReturn}. None of them is a commit. An invocation that is
canceled after Start never completes: its task may instead be abandoned
(\rulename{WORKER-ABANDON}), before or after any of its effects, or
dropped as stale once it has returned (\rulename{STEP-COMPLETE-STALE}).
Any other enabled transition can occur between any two of these. The
method
reads the snapshot taken at Start. Its writes are applied blindly at
Complete; nothing checks that the values it read are still current. So a
Step method is \emph{not} a read-modify-write transaction: two Steps that
both read an Attribute and both write it can lose an update. The
translation is designed around this: it puts at most one shared access in
each method, so that a read linearizes at Start and a write at Complete.

\paragraph{Inside one commit.}
The effect of a commit is defined by deterministic meta-functions
(\file{Dex/MetaFunctions.lean}), run in a fixed order. The ones the proof
relies on are these.
\begin{itemize}
\item \emph{Data effects}: Attribute writes in order, then Channel
  publications in order, then a \emph{wake} of the Channels published to.
\item \emph{Wake}: every Step that waits on a changed Channel is
  re-evaluated, earliest registration first (by $\mathit{sws}$). With
  \kw{ForOne} waiters only, a Channel is therefore FIFO by registration.
  Lean: \lean{Dex.wake_servesEarliestWaiter}.
\item \emph{Evaluation}: a wait is satisfied when its plan is feasible.
  \kw{AllOf} needs every condition. \kw{AnyOf} takes the first feasible
  candidate in the order timers, then Channel conditions in declaration
  order, and consumes only from the winner
  (\lean{Dex.registerWait_any_firstChannelWins}). A satisfied wait moves the
  Step to phase $\mathit{EX}$ and makes its \kw{Execute} pending.
\item \emph{Registration}: a non-empty wait arms its timers at the commit's
  timestamp plus whole seconds, records $\mathit{sws}$, and is evaluated
  once. A timer is never done when its wait registers, even a zero-second
  one (\lean{Dex.registerWait_timer_blocks}).
\item \emph{Decision}: the deciding Step completes; then a close decision
  closes the run, canceling every active Step, or the moves create new Step
  executions.
\end{itemize}

\paragraph{Timers fire late and in order.}
\rulename{TIMER-FIRE} is enabled when some armed timer's fire time is at or
before the clock. It fires a non-empty \emph{prefix} of the due timers, in
the order (fire time, encoded timer ID), in one commit. Firing can be
arbitrarily late: any commit may come between a timer becoming due and its
firing, so a due timer can lose an \kw{AnyOf} race to a message that arrives
after its fire time. The specification argues from the timer service's code
that the server fires a prefix, and not an arbitrary subset. A later review
found one window, without any failure, in which that is not exact: the
server can fire a timer while an earlier-due timer, armed by a command
applied just before, is still armed. Both theorems keep the prefix rule as
an explicit idealization (\S\ref{sec:limitations}). An over-approximation
would be unsafe here: for an equivalence theorem, extra Dex runs are extra
obligations for \Lang{}.

\subsection{Observable events}
\label{sec:dex-obs}

The Dex event alphabet \Obs{} of the restricted system is
$\ev{start}(T, v)$, $\ev{pub}(c, v)$, $\ev{rejected}(\ev{pub}(c, v))$,
$\ev{fx}(e, m, a)$, and $\ev{closed}(\mathit{status}, \mathit{outs})$
(\lean{Dex.Obs}). An $\ev{fx}$ event is one external effect of a handler,
such as a call to another system. It is emitted by the Worker, outside any
commit: an effect can happen with no commit after it, and after the Flow
closed, when its Step was canceled in between.

Reads are not observable. Including them would make every intermediate
Attribute write visible and would turn the claim into a much stronger one
that no compilation of threads into Steps could satisfy.
